% =====================================================================
\chapter{Cuantificadores y negación}
\label{tema:2}
% =====================================================================

Casi ningún teorema habla de un único objeto: afirman algo de \emph{todos} los enteros, de \emph{alguna} raíz de un polinomio o de que \emph{para cada} $\varepsilon$ \emph{existe} un $\delta$. Este tema introduce los predicados y los cuantificadores, estudia el efecto del orden de los cuantificadores y da reglas sistemáticas para negar enunciados, que serán imprescindibles en las demostraciones por contrarrecíproco y por reducción al absurdo. Se siguen las §2.7–2.10 de~\cite{hammack}, el capítulo~2 de~\cite{velleman} (§2.1 y §2.2), la §3.6 de~\cite{mcs} y la guía de negación del curso CS103~\cite{cs103}.

\section{Predicados}

\begin{definicion}[Predicado]\label{def:predicado}
Sea $U$ un conjunto no vacío, llamado \emph{universo} o dominio. Un \emph{predicado} sobre $U$ es un enunciado $P(x)$ que contiene una variable $x$ y que se convierte en una proposición cada vez que se sustituye $x$ por un elemento concreto de $U$. El \emph{conjunto de verdad} de $P$ es $\{x\in U: P(x)\text{ es verdadera}\}$. Del mismo modo se definen predicados de varias variables, $P(x,y)$, con $(x,y)\in U\times W$.
\end{definicion}

\begin{ejemplo}
Con $U=\Z$, el predicado $P(n)$: «$n^{2}-1$ es múltiplo de~8» da proposiciones como $P(3)$ ($8\divides 8$, verdadera) o $P(2)$ ($8\nodivide 3$, falsa). Su conjunto de verdad es el de los enteros impares (\cref{eje:4-24} prueba una versión más fuerte).
\end{ejemplo}

\section{Cuantificador universal y existencial}

\begin{definicion}[Cuantificadores]\label{def:cuantificadores}
Sea $P(x)$ un predicado sobre $U$.
\begin{enumerate}
  \item $\forall x\in U\ P(x)$ («para todo $x$ de $U$, $P(x)$») es verdadera si $P(a)$ es verdadera para cada $a\in U$.
  \item $\exists x\in U\ P(x)$ («existe $x$ en $U$ tal que $P(x)$») es verdadera si $P(a)$ es verdadera para al menos un $a\in U$.
\end{enumerate}
Cuando el universo está claro se escribe $\forall x\,P(x)$ y $\exists x\,P(x)$. La variable $x$ está \emph{ligada} por el cuantificador: $\forall x\,P(x)$ y $\forall t\,P(t)$ dicen lo mismo.
\end{definicion}

Los cuantificadores se combinan con las conectivas. Las formas más frecuentes son la \emph{universal condicional}, $\forall x\,\bigl(P(x)\Rightarrow Q(x)\bigr)$, que traduce «todo $P$ es $Q$», y la \emph{existencial conjuntiva}, $\exists x\,\bigl(P(x)\wedge Q(x)\bigr)$, que traduce «algún $P$ es $Q$». Se abrevian como $\forall x\in A\ Q(x)$ y $\exists x\in A\ Q(x)$ cuando $P(x)$ es «$x\in A$».

\begin{ejemplo}\label{ej:traduccion}
En $U=\Z$: «todo cuadrado perfecto es no negativo» es $\forall n\,(n^{2}\ge 0)$; «hay un primo par» es $\exists p\,(p\text{ primo}\wedge p\text{ par})$; «ningún entero es par e impar a la vez» es $\neg\exists n\,(n\text{ par}\wedge n\text{ impar})$, o bien $\forall n\,\neg(n\text{ par}\wedge n\text{ impar})$.
\end{ejemplo}

\begin{observacion}
Si $A=\varnothing$, el enunciado $\forall x\in A\ Q(x)$ es verdadero, sea cual sea $Q$: no hay ningún $x\in A$ que lo incumpla. Es la versión cuantificada del condicional vacuamente verdadero del \cref{tema:1}. En cambio, $\exists x\in\varnothing\ Q(x)$ es siempre falso.
\end{observacion}

\section{Cuantificadores anidados}

Con predicados de varias variables, los cuantificadores se anidan. El orden importa cuando se mezclan cuantificadores de distinto tipo.

\begin{ejemplo}\label{ej:orden}
Sea $P(x,y)$: «$x+y=0$» con $x,y\in\Z$. El enunciado $\forall x\,\exists y\,P(x,y)$ dice que todo entero tiene un opuesto, y es verdadero: dado $x$, sirve $y=-x$. El enunciado $\exists y\,\forall x\,P(x,y)$ dice que hay un entero $y$ que es opuesto de \emph{todos} los enteros a la vez, y es falso: si existiera, $0+y=0$ y $1+y=0$ darían $y=0$ e $y=-1$. En el primer enunciado $y$ puede depender de $x$; en el segundo, no.
\end{ejemplo}

\begin{proposicion}\label{prop:orden}
Para todo predicado $P(x,y)$,
\[
\exists y\,\forall x\,P(x,y)\ \Longrightarrow\ \forall x\,\exists y\,P(x,y),
\]
y la implicación recíproca es falsa en general.
\end{proposicion}

\begin{proof}
Supóngase $\exists y\,\forall x\,P(x,y)$ y sea $y_0$ un elemento tal que $P(x,y_0)$ es verdadera para todo $x$. Sea $x$ arbitrario. Entonces $P(x,y_0)$ es verdadera, y por tanto existe $y$ (a saber, $y_0$) con $P(x,y)$. Como $x$ era arbitrario, $\forall x\,\exists y\,P(x,y)$. El \cref{ej:orden} muestra que la recíproca puede fallar.
\end{proof}

Cuantificadores del mismo tipo sí pueden intercambiarse: $\forall x\,\forall y\,P(x,y)\eqlog\forall y\,\forall x\,P(x,y)$ y $\exists x\,\exists y\,P(x,y)\eqlog\exists y\,\exists x\,P(x,y)$.

\begin{ejemplo}[Convergencia de una sucesión]\label{ej:limite}
Que una sucesión real $(a_n)_{n\in\N}$ converja a $L$ se escribe
\[
\forall\varepsilon>0\ \exists N\in\N\ \forall n\in\N\ \bigl(n\ge N\Rightarrow\abs{a_n-L}<\varepsilon\bigr).
\]
El índice $N$ puede depender de $\varepsilon$: cuanto más exigente es la tolerancia, más lejos hay que ir en la sucesión. Intercambiar los dos primeros cuantificadores daría un enunciado mucho más fuerte y casi nunca cierto (\cref{eje:2-arquimedes}).
\end{ejemplo}

\section{Negación de enunciados cuantificados}

\begin{teorema}[Negación de cuantificadores]\label{teo:negacion}
Para todo predicado $P(x)$ sobre un universo $U$,
\[
\neg\,\forall x\,P(x)\ \eqlog\ \exists x\,\neg P(x),\qquad
\neg\,\exists x\,P(x)\ \eqlog\ \forall x\,\neg P(x).
\]
\end{teorema}

\begin{proof}
$\neg\,\forall x\,P(x)$ es verdadera si y solo si no es cierto que $P(a)$ valga para todo $a\in U$, es decir, si y solo si hay algún $a\in U$ con $P(a)$ falsa, que es la condición para que $\exists x\,\neg P(x)$ sea verdadera. Para la segunda equivalencia se aplica la primera al predicado $\neg P$ y se niegan ambos miembros: $\neg\,\forall x\,\neg P(x)\eqlog\exists x\,\neg\neg P(x)\eqlog\exists x\,P(x)$; negando, $\forall x\,\neg P(x)\eqlog\neg\,\exists x\,P(x)$.
\end{proof}

El \cref{teo:negacion} y las leyes de la \cref{prop:leyes} dan un procedimiento mecánico: la negación «atraviesa» los cuantificadores cambiando $\forall$ por $\exists$ y viceversa, y al llegar a la parte proposicional se aplican De Morgan y la negación del condicional, $\neg(P\Rightarrow Q)\eqlog P\wedge\neg Q$. Con universos acotados, el cuantificador conserva su dominio: $\neg\,\forall x\in A\ Q(x)\eqlog\exists x\in A\ \neg Q(x)$.

\begin{ejemplo}\label{ej:negar-limite}
Se niega la definición de convergencia del \cref{ej:limite}:
\begin{align*}
&\neg\,\forall\varepsilon>0\ \exists N\ \forall n\ \bigl(n\ge N\Rightarrow\abs{a_n-L}<\varepsilon\bigr)\\
\eqlog\ &\exists\varepsilon>0\ \forall N\ \exists n\ \neg\bigl(n\ge N\Rightarrow\abs{a_n-L}<\varepsilon\bigr)\\
\eqlog\ &\exists\varepsilon>0\ \forall N\ \exists n\ \bigl(n\ge N\wedge\abs{a_n-L}\ge\varepsilon\bigr).
\end{align*}
Es decir, $(a_n)$ \emph{no} converge a $L$ si hay una tolerancia $\varepsilon$ tal que, por lejos que se mire, siempre aparece algún término a distancia al menos $\varepsilon$ de $L$. Con $a_n=(-1)^{n}$ y $L=1$ sirve $\varepsilon=1$: para cualquier $N$, el término impar $n=2N+1\ge N$ cumple $\abs{a_n-1}=2\ge 1$.
\end{ejemplo}

\section{Traducción y unicidad}

Las expresiones del lenguaje natural ocultan a menudo los cuantificadores. «Los primos mayores que~2 son impares» es universal; «un cuadrado puede ser par» es existencial; «solo los múltiplos de~6 son múltiplos de~2 y de~3» significa «si $n$ es múltiplo de~2 y de~3, entonces es múltiplo de~6». Antes de demostrar un enunciado conviene escribirlo con cuantificadores explícitos~\cite[§2.9]{hammack}.

\begin{definicion}[Existencia y unicidad]\label{def:unicidad}
$\exists!\,x\,P(x)$ («existe un único $x$ tal que $P(x)$») abrevia
\[
\exists x\,\Bigl(P(x)\wedge\forall y\,\bigl(P(y)\Rightarrow y=x\bigr)\Bigr).
\]
\end{definicion}

Así, $\exists!$ contiene dos afirmaciones: hay al menos un elemento con la propiedad (\emph{existencia}) y dos elementos con la propiedad son iguales (\emph{unicidad}). El \cref{tema:6} explica cómo se demuestra cada parte.

\begin{observacion}[Errores frecuentes]\label{obs:2-errores}
(a) Negar «todos los $x$ cumplen $P$» como «ningún $x$ cumple $P$»: la negación correcta es «algún $x$ no cumple $P$».
(b) Negar «$\forall x\,(P(x)\Rightarrow Q(x))$» como otro condicional: la negación es $\exists x\,(P(x)\wedge\neg Q(x))$, una conjunción.
(c) Cambiar el orden de $\forall$ y $\exists$ al copiar una definición (\cref{ej:orden}).
(d) Olvidar el universo: $\exists x\,(x^{2}=2)$ es verdadero en $\R$ y falso en $\Q$.
\end{observacion}

\begin{observacion}[Conexión con el curso]
La definición $\varepsilon$–$N$ del \cref{ej:limite} y su análoga $\varepsilon$–$\delta$ para funciones son la base de los límites y derivadas del módulo~4, de los resultados de convergencia de los algoritmos de optimización del módulo~7 y del análisis de la softmax con temperatura del módulo~1, donde se estudia el límite cuando la temperatura tiende a~0. Saber negar estas definiciones es lo que permite demostrar que un algoritmo \emph{no} converge o que una función \emph{no} es continua.
\end{observacion}

\begin{lecturas}
\item Hammack~\cite{hammack}, §2.7 «Quantifiers», §2.9 «Translating English to Symbolic Logic» y §2.10 «Negating Statements».
\item Velleman~\cite{velleman}, §2.1 «Quantifiers» y §2.2 «Equivalences Involving Quantifiers».
\item Lehman, Leighton y Meyer~\cite{mcs}, §1.2 «Predicates» y §3.6 «Predicate Formulas».
\item Stanford CS103~\cite{cs103}, clases~4 y~5 «First-Order Logic» y la guía «Guide to Negation».
\end{lecturas}

\section*{Ejercicios del tema 2}
\addcontentsline{toc}{section}{Ejercicios del tema 2}

\begin{ejercicio}{1}{eje:2-valores}
Determine el valor de verdad: (a) $\forall x\in\R\ (x^{2}\ge 0)$; (b) $\exists x\in\Q\ (x^{2}=2)$; (c) $\exists x\in\R\ (x^{2}=2)$; (d) $\forall n\in\N\ \exists m\in\N\ (m<n)$; (e) $\exists m\in\N\ \forall n\in\N\ (m\le n)$; (f) $\forall x\in\varnothing\ (x\neq x)$.
\end{ejercicio}

\begin{ejercicio}{1}{eje:2-simbolos}
Escriba con cuantificadores, sobre el universo indicado: (a) «todo entero par mayor que~2 es suma de dos primos» ($\Z$); (b) «hay un número real que es mayor que todos los enteros» y decida si es verdadero; (c) «ningún cuadrado de un entero termina en~2»; (d) «la ecuación $x^{2}+1=0$ no tiene solución real».
\end{ejercicio}

\begin{ejercicio}{1}{eje:2-lenguaje}
Exprese en lenguaje natural y determine el valor de verdad: (a) $\forall x\in\R\ \exists y\in\R\ (y>x)$; (b) $\exists y\in\R\ \forall x\in\R\ (y>x)$; (c) $\forall x\in\R\ \bigl(x\neq 0\Rightarrow\exists y\in\R\ (xy=1)\bigr)$; (d) $\exists x\in\Z\ \forall y\in\Z\ (xy=y)$.
\end{ejercicio}

\begin{ejercicio}{2}{eje:2-negar}
Niegue los enunciados siguientes de modo que el símbolo $\neg$ no aparezca delante de ningún cuantificador ni paréntesis, y decida cuál es verdadero, el original o su negación: (a) $\forall x\in\R\ \exists y\in\R\ (x<y<x+1)$; (b) $\exists n\in\N\ \forall m\in\N\ (n\divides m)$; (c) $\forall x\in\R\ \bigl(x^{2}>4\Rightarrow x>2\bigr)$; (d) $\forall n\in\Z\ \bigl(n\text{ primo}\Rightarrow n\text{ impar}\vee n=2\bigr)$.
\fuente{adaptado de Hammack~\cite{hammack}, §2.10 «Negating Statements».}
\end{ejercicio}

\begin{ejercicio}{2}{eje:2-continuidad}
Una función $f\colon\R\to\R$ es continua en $a$ si $\forall\varepsilon>0\ \exists\delta>0\ \forall x\in\R\ \bigl(\abs{x-a}<\delta\Rightarrow\abs{f(x)-f(a)}<\varepsilon\bigr)$. (a) Escriba la negación sin $\neg$ delante de cuantificadores. (b) Use la negación para demostrar que la función $f(x)=0$ si $x<0$ y $f(x)=1$ si $x\ge 0$ no es continua en $a=0$.
\fuente{adaptado de Velleman~\cite{velleman}, §2.2 «Equivalences Involving Quantifiers».}
\end{ejercicio}

\begin{ejercicio}{2}{eje:2-orden}
Para cada predicado, determine el valor de verdad de $\forall x\,\exists y\,P(x,y)$ y de $\exists y\,\forall x\,P(x,y)$: (a) $P(x,y)$: «$x+y=0$» en $\N_0$; (b) $P(x,y)$: «$xy=1$» en $\Q\setminus\{0\}$; (c) $P(x,y)$: «$x\le y$» en $\N$; (d) $P(x,y)$: «$y\le x$» en $\N$.
\end{ejercicio}

\begin{ejercicio}{2}{eje:2-distribuir}
Decida cuáles de las siguientes equivalencias son ciertas para todos los predicados $P$ y $Q$ sobre cualquier universo, y dé un contraejemplo para las que no lo sean: (a) $\forall x\,(P(x)\wedge Q(x))\eqlog\forall x\,P(x)\wedge\forall x\,Q(x)$; (b) $\forall x\,(P(x)\vee Q(x))\eqlog\forall x\,P(x)\vee\forall x\,Q(x)$; (c) $\exists x\,(P(x)\vee Q(x))\eqlog\exists x\,P(x)\vee\exists x\,Q(x)$; (d) $\exists x\,(P(x)\wedge Q(x))\eqlog\exists x\,P(x)\wedge\exists x\,Q(x)$.
\fuente{adaptado de Velleman~\cite{velleman}, §2.2 «Equivalences Involving Quantifiers».}
\end{ejercicio}

\begin{ejercicio}{3}{eje:2-distribuir-dem}
Demuestre las equivalencias (a) y (c) del \cref{eje:2-distribuir}. Demuestre además que $\forall x\,P(x)\vee\forall x\,Q(x)\Rightarrow\forall x\,(P(x)\vee Q(x))$ es siempre verdadera.
\end{ejercicio}

\begin{ejercicio}{3}{eje:2-vacio}
Demuestre que, si el universo $U$ es no vacío, $\forall x\,P(x)\Rightarrow\exists x\,P(x)$ para todo predicado $P$. Explique por qué la implicación falla si se permite $U=\varnothing$.
\end{ejercicio}

\begin{ejercicio}{2}{eje:2-unicidad}
(a) Escriba $\exists!\,x\,P(x)$ sin el símbolo $\exists!$ de una forma distinta de la \cref{def:unicidad}, como conjunción de un enunciado de existencia y otro de unicidad. (b) Niegue $\exists!\,x\,P(x)$ y explique en palabras las dos maneras en que puede fallar. (c) Decida si $\exists!\,x\in\R\ (x^{2}=4)$ y $\exists!\,x\in\N\ (x^{2}=4)$ son verdaderos.
\end{ejercicio}

\begin{ejercicio}{2}{eje:2-acotada}
Una función $f\colon\R\to\R$ es \emph{acotada} si $\exists M\in\R\ \forall x\in\R\ (\abs{f(x)}\le M)$, y es \emph{creciente} si $\forall x\,\forall y\ (x<y\Rightarrow f(x)\le f(y))$. (a) Niegue ambas definiciones. (b) Demuestre, usando la negación, que $f(x)=x^{2}$ no es acotada ni creciente.
\end{ejercicio}

\begin{ejercicio}{3}{eje:2-arquimedes}
Se admite la \emph{propiedad arquimediana}: para todo número real $x$ existe $N\in\N$ con $N>x$. Demuestre que (a) $\forall\varepsilon>0\ \exists N\in\N\ \forall n\in\N\ (n\ge N\Rightarrow 1/n<\varepsilon)$ es verdadero, y que (b) $\exists N\in\N\ \forall\varepsilon>0\ \forall n\in\N\ (n\ge N\Rightarrow 1/n<\varepsilon)$ es falso.
\end{ejercicio}

\begin{ejercicio}{2}{eje:2-critica}
Se pretende demostrar que «todo número natural es impar» así: «Por reducción al absurdo, supóngase que todo número natural es par. Entonces $137$ es par, pero $137=2\cdot 68+1$ es impar. Contradicción; luego todo natural es impar.» Identifique el error lógico, escriba la negación correcta del enunciado y decida si este es verdadero.
\fuente{adaptado de Stanford CS103~\cite{cs103}, \emph{Problem Set}~1, problema~9 «Proof Critiques».}
\end{ejercicio}

\begin{ejercicio}{3}{eje:2-finito}
Sea $U=\{a_1,\dots,a_n\}$ un universo finito. Justifique que $\forall x\,P(x)\eqlog P(a_1)\wedge\dots\wedge P(a_n)$ y que $\exists x\,P(x)\eqlog P(a_1)\vee\dots\vee P(a_n)$, y deduzca el \cref{teo:negacion} para universos finitos a partir de las leyes de De Morgan generalizadas del \cref{eje:1-demorgan-n}.
\end{ejercicio}
